%! AP Statistics — Unit 7, Lesson 7.1 — Guided Notes (Answer Key)
\documentclass[10pt]{article}
\usepackage{apstats-article}
\usepackage{apstats-key}

\begin{document}

\pageheader{Unit 7, Lesson 7.1}{Guided Notes --- Answer Key}
\namedateperiod

\begin{objectivebox}
By the end of this lesson, I will be able to:
\begin{itemize}
    \item[$\square$] \ansline{Describe and identify Type I and Type II errors in context.}
    \item[$\square$] \ansline{Explain how $\alpha$ controls the Type I error rate and affects Type II error.}
    \item[$\square$] \ansline{Explain why inference for means uses the $t$-distribution instead of $z$.}
\end{itemize}
\end{objectivebox}

\begin{vocabbox}
\termblanklong{Type I error}
\ansline{Rejecting a true null hypothesis (false positive); probability $= \alpha$.}

\termblanklong{Type II error}
\ansline{Failing to reject a false null hypothesis (false negative); probability $= \beta$.}

\termblanklong{Significance level ($\alpha$)}
\ansline{The maximum tolerable probability of a Type I error; set before data collection.}

\termblanklong{$t$-distribution}
\ansline{A family of bell-shaped, symmetric distributions with heavier tails than the normal; used when $\sigma$ is unknown and $s$ is substituted.}
\end{vocabbox}

\begin{hookbox}
\textbf{Discussion:} A drug trial finds a statistically significant result. A follow-up
study contradicts it. How is this possible?

\begin{teachernote}
    Both studies may have been conducted correctly. If the initial study committed a Type I
    error (rejected a true $H_0$), a follow-up study using a different sample is likely to
    fail to reject $H_0$. This is why replication is essential in science.
\end{teachernote}
\end{hookbox}

\begin{notesbox}{Errors in Statistical Inference}
An inference procedure can produce two types of incorrect decisions:

\renewcommand{\arraystretch}{2}
\begin{tabularx}{\linewidth}{lXX}
    & \textbf{$H_0$ is True} & \textbf{$H_0$ is False} \\
    \hline
    \textbf{Reject $H_0$} & \ans{Type I error} & Correct decision \\
    \textbf{Fail to reject $H_0$} & Correct decision & \ans{Type II error} \\
\end{tabularx}

\vspace{0.2in}
\textbf{Type I error:} rejecting a \ans{true} null hypothesis. \\[4pt]
The probability of a Type I error $=$ \ans{$\alpha$} (the significance level).

\vspace{0.15in}
\textbf{Type II error:} failing to reject a \ans{false} null hypothesis. \\[4pt]
The probability of a Type II error $= \beta$ (beta). \textbf{Power} $= 1 - \beta$.
\end{notesbox}

\begin{notesbox}{The Trade-off Between Error Types}
\textbf{Key Principle:} For a fixed sample size, decreasing the probability of one type
of error \ans{increases} the probability of the other.

\vspace{0.1in}
If we lower $\alpha$ from 0.05 to 0.01:
\begin{itemize}
    \item Type I error rate \ans{decreases}
    \item Type II error rate \ans{increases}
    \item Power \ans{decreases}
\end{itemize}

\vspace{0.1in}
To reduce \emph{both} error rates simultaneously, we must \ans{increase the sample size $n$}.
\end{notesbox}

\begin{notesbox}{Connecting to Unit 7 --- Why the $t$-Distribution?}
In Unit 6, inference for proportions used the \ans{normal ($z$)} distribution because
$\sigma_{\hat p} = \sqrt{p(1-p)/n}$ could be estimated from the hypothesized $p$.

For means, $\sigma$ is \ans{rarely known in practice}, so we substitute the sample standard
deviation $s$. This introduces extra variability, producing the \ans{$t$}-distribution.

\vspace{0.1in}
The $t$-statistic: \quad $t = \dfrac{\bar x - \mu_0}{s/\sqrt{n}}$ \quad follows a
$t$-distribution with $df = $ \ans{$n - 1$} degrees of freedom.

\vspace{0.1in}
The $t$-distribution has \ans{heavier} tails than the standard normal.
As $df$ increases, the $t$-distribution approaches the \ans{standard normal ($z$)} distribution.
\end{notesbox}

\begin{practicebox}
\textbf{Guided Practice:} A hospital tests whether a new pain medication reduces
average recovery time below 14 days ($H_0{:}\ \mu = 14$, $H_a{:}\ \mu < 14$).

\begin{enumerate}
  \item Describe a Type I error in context.
        \ansline{Concluding the medication reduces mean recovery time below 14 days when it actually does not.}

  \item Describe a Type II error in context.
        \ansline{Failing to conclude the medication reduces mean recovery time when it actually does reduce it below 14 days.}

  \item Which error is more costly for patients? Explain.
        \ansline{Answers vary. A Type I error means giving patients an ineffective drug; a Type II error means withholding a beneficial one. Both are costly; the answer depends on context and risk tolerance.}

  \item If $\alpha$ is lowered from 0.05 to 0.01, describe the effect on both
        error types.
        \ansline{Type I error rate decreases (from 5\% to 1\%). Type II error rate increases (we require stronger evidence to reject, so we are more likely to miss a real effect). Power decreases.}
\end{enumerate}
\end{practicebox}

\end{document}
